\textbf{Protocolo de Validación Científica (TUI v4.2 - Fase 2)}

Título del Experimento: Análisis de la Tensión: Verificación de la
Hipótesis H1 bajo risk\_scale variable.

Alias: "Resolución del Problema 3 (P3)"

Investigador Principal: (Jose M Rivera Garcia)

\textbf{1. Definición del Problema de Investigación}

El proyecto TUI v4.1 se enfrenta a una aparente contradicción que debe
ser resuelta:

\begin{enumerate}
\def\labelenumi{\arabic{enumi}.}
\item
  \textbf{La Teoría (H1):} Nuestra tesis central postula que la
  inteligencia operativa emerge como función del riesgo acumulado (\$I
  \textbackslash propto P\_\{riesgo\}\$). La implicación es que un
  entorno con mayor riesgo debería fomentar un mayor desarrollo de
  inteligencia prudencial.
\item
  \textbf{El Dato Observado (P3):} Los resultados preliminares del
  "barrido de riesgo" (sweep\_risk.json) muestran lo contrario. A medida
  que el risk\_scale (el riesgo ambiental) aumenta, la avg\_reward
  (recompensa neta) del agente Simbiosis \emph{disminuye} (de +588.19 a
  +556.91).
\end{enumerate}

Pregunta Central de Investigación:

¿Por qué la recompensa neta del agente (nuestro proxy de éxito)
disminuye cuando el riesgo ambiental aumenta, si nuestra teoría H1
predice que el riesgo fomenta la inteligencia? ¿Es este dato una
refutación de la H1?

\textbf{2. El Hallazgo (Hipótesis de Trabajo)}

Nuestra hipótesis de trabajo es que estos datos \textbf{no refutan} la
H1. Por el contrario, la \emph{refuerzan} al revelar una tensión
dinámica que demuestra la naturaleza \emph{prudencial} del agente PGF.

El error está en equiparar la avg\_reward neta (un puntaje final) con la
"inteligencia" (la calidad del proceso de decisión).

Proponemos que la avg\_reward (el PGF Neto) debe desacoplarse en dos
componentes, tal como se define en la lógica de evaluator\_pgf.py (antes
en dqn\_agent.py):

PGF\_Neto = PGF\_Beneficio\_Bruto - PGF\_Costo\_Ambiental

\begin{enumerate}
\def\labelenumi{\arabic{enumi}.}
\item
  \textbf{PGF\_Beneficio\_Bruto:} Es el término kappa * delta\_P * A\_t.
  Esta es nuestra verdadera métrica de "inteligencia": la habilidad del
  agente para tomar decisiones que reducen el riesgo futuro y se alinean
  con el propósito.
\item
  \textbf{PGF\_Costo\_Ambiental:} Es el término lambda\_c * delta\_C\_t.
  Este término mide las penalizaciones de recursos sufridas en el
  entorno (ej. pisar \emph{tripwires}).
\end{enumerate}

El script prototipo\_rl\_simbiosis.py multiplica las penalizaciones por
el risk\_scale. Por lo tanto, nuestra hipótesis es:

\textbf{"Al aumentar el risk\_scale, el PGF\_Costo\_Ambiental aumenta
drásticamente (los errores son más caros). El agente se vuelve
\emph{más} inteligente y prudente (su PGF\_Beneficio\_Bruto se mantiene
alto), pero su puntuación PGF\_Neto disminuye porque los costos del
entorno son mayores. La caída en la recompensa neta no es un fracaso de
inteligencia; es la evidencia de una gestión de riesgos exitosa en un
entorno hostil."}

\textbf{3. Analogías para el Análisis}

Para comunicar este hallazgo, usaremos dos analogías:

\begin{enumerate}
\def\labelenumi{\arabic{enumi}.}
\item
  \textbf{Analogía del Piloto de F1:}

  \begin{itemize}
  \item
    \textbf{Agente:} Un piloto de F1 de clase mundial (inteligencia
    alta).
  \item
    \textbf{risk\_scale:} El clima. risk\_scale=0.5 es un día soleado.
    risk\_scale=2.0 es un aguacero torrencial.
  \item
    \textbf{PGF\_Beneficio\_Bruto:} La \emph{habilidad} del piloto para
    trazar curvas y controlar el coche. Esta habilidad es alta en ambos
    climas.
  \item
    \textbf{PGF\_Costo\_Ambiental:} El costo de los errores. En el sol,
    un pequeño derrape cuesta 0.1 seg. En la lluvia, el \emph{mismo}
    derrape cuesta 5 seg.
  \item
    \textbf{PGF\_Neto (Resultado):} El tiempo de vuelta final. El piloto
    es un genio en la lluvia (PGF Bruto alto), pero su tiempo de vuelta
    (PGF Neto) será inevitablemente más lento que en el día soleado,
    porque el entorno castiga más duramente los errores.
  \end{itemize}
\item
  \textbf{Analogía del CEO en Recesión:}

  \begin{itemize}
  \item
    \textbf{Agente:} Un CEO inteligente.
  \item
    \textbf{risk\_scale:} La economía. risk\_scale=0.5 es un mercado en
    auge. risk\_scale=2.0 es una recesión severa.
  \item
    \textbf{PGF\_Beneficio\_Bruto:} La \emph{calidad} de las decisiones
    estratégicas del CEO.
  \item
    \textbf{PGF\_Costo\_Ambiental:} El costo de operación (intereses,
    despidos, etc.).
  \item
    \textbf{PGF\_Neto (Resultado):} Las ganancias netas de la empresa.
    En una recesión (risk\_scale alto), un CEO inteligente (PGF Bruto
    alto) no busca ganancias récord; busca \emph{minimizar pérdidas},
    \emph{cortar costos} y \emph{sobrevivir}. Sus ganancias netas (PGF
    Neto) serán bajas, pero esto es un acto de inteligencia superior
    comparado con un CEO ingenuo que quiebra por ser demasiado agresivo.
  \end{itemize}
\end{enumerate}

\textbf{4. Requerimientos del Simulador (Modificaciones v4.2)}

El simulador, tras la refactorización de la Fase 1
(feature/tui-v4.2-refactorizacion-metodologica), ya está
metodológicamente listo. Sin embargo, para probar esta hipótesis (P3),
requerimos una modificación en el \emph{logging}:

\begin{enumerate}
\def\labelenumi{\arabic{enumi}.}
\item
  \textbf{Modificar evaluator\_pgf.py:}

  \begin{itemize}
  \item
    La función calcular\_metricas (o como se llame ahora) actualmente
    calcula y retorna un solo valor (self.PGF).
  \item
    Debe ser modificada para calcular y retornar \emph{tres} valores por
    separado:

    \begin{enumerate}
    \def\labelenumii{\arabic{enumii}.}
    \item
      pgf\_beneficio\_bruto = (kappa * delta\_P * A\_t)
    \item
      pgf\_costo\_ambiental = (lambda\_c * delta\_C\_t)
    \item
      pgf\_neto = pgf\_beneficio\_bruto - pgf\_costo\_ambiental
    \end{enumerate}
  \end{itemize}
\item
  \textbf{Modificar prototipo\_rl\_simbiosis.py:}

  \begin{itemize}
  \item
    El bucle de entrenamiento debe ahora recibir estos tres valores del
    evaluador.
  \item
    Debe registrar el pgf\_neto para el aprendizaje del agente
    (agent.learn(..., reward=pgf\_neto, ...)).
  \item
    Debe registrar los tres valores en el log del episodio.
  \item
    La exportación final a .json y .csv debe incluir las medias de estos
    tres componentes por separado (ej. avg\_pgf\_neto,
    avg\_pgf\_beneficio\_bruto, avg\_pgf\_costo\_bruto).
  \end{itemize}
\end{enumerate}

\textbf{5. Procedimiento Científico de Validación (Fase 2)}

Para que el procedimiento sea científicamente correcto, debe ser
riguroso, comparativo y falsable.

\begin{enumerate}
\def\labelenumi{\arabic{enumi}.}
\item
  \textbf{Configuración Experimental:}

  \begin{itemize}
  \item
    \textbf{Agentes:} Usar los 3 agentes definidos en Fase 1 (Q-Control,
    DQN-Control, DQN-Simbiosis).
  \item
    \textbf{Semillas:} Ejecutar cada condición con un mínimo de 5
    semillas (ej. -\/-seeds 42, 123, 456, 789, 1000) para asegurar la
    reproducibilidad y obtener intervalos de confianza.
  \item
    \textbf{Barrido (Sweep):} Ejecutar el barrido de risk\_scale en al
    menos 4 niveles (ej. {[}0.5, 1.0, 1.5, 2.0{]}).
  \end{itemize}
\item
  \textbf{Recolección de Datos:}

  \begin{itemize}
  \item
    Ejecutar el sweep completo (3 agentes x 5 semillas x 4 niveles de
    riesgo = 60 corridas de 1000 episodios).
  \item
    Agregar los resultados en un archivo .csv o .json que contenga las
    medias por episodio para: risk\_scale, seed, agent\_type,
    avg\_pgf\_neto, avg\_pgf\_beneficio\_bruto, avg\_pgf\_costo\_bruto.
  \end{itemize}
\item
  \textbf{Análisis y Visualización:}

  \begin{itemize}
  \item
    \textbf{Gráfico 1 (Confirmar P3):} Graficar risk\_scale (Eje X) vs.
    avg\_pgf\_neto (Eje Y).

    \begin{itemize}
    \item
      \emph{Predicción:} Para el agente Simbiosis, esta línea será
      descendente, confirmando el dato observado.
    \end{itemize}
  \item
    \textbf{Gráfico 2 (La Prueba de H1):} Graficar risk\_scale (Eje X)
    vs. avg\_pgf\_beneficio\_bruto (Eje Y).

    \begin{itemize}
    \item
      \emph{Predicción:} Esta línea será \textbf{plana o ascendente},
      mostrando que la "inteligencia" (habilidad prudencial) del agente
      no disminuye.
    \end{itemize}
  \item
    \textbf{Gráfico 3 (El Costo):} Graficar risk\_scale (Eje X) vs.
    avg\_pgf\_costo\_bruto (Eje Y).

    \begin{itemize}
    \item
      \emph{Predicción:} Esta línea será \textbf{fuertemente
      ascendente}, mostrando que el costo ambiental es el único
      responsable de la caída en el Gráfico 1.
    \end{itemize}
  \end{itemize}
\item
  \textbf{Criterios de Éxito y Falsabilidad (Conclusión Científica):}

  \begin{itemize}
  \item
    \textbf{Éxito (Hipótesis Confirmada):} Si el Gráfico 1 (Neto)
    desciende, PERO el Gráfico 2 (Bruto) se mantiene plano o asciende.

    \begin{itemize}
    \item
      \emph{Conclusión del Paper:} "Demostramos que el PGF crea un
      agente prudencial. A medida que el riesgo ambiental aumenta, la
      recompensa neta del agente disminuye, no por una falla de
      inteligencia, sino porque los costos de error se magnifican. El
      agente mantiene o incrementa su generación de
      \textquotesingle beneficio PGF\textquotesingle{} (Gráfico 2)
      mientras gestiona costos mayores (Gráfico 3), validando la H1 a un
      nivel más profundo."
    \end{itemize}
  \item
    \textbf{Fallo (Hipótesis Falsada):} Si el Gráfico 2 (Beneficio
    Bruto) también desciende bruscamente junto con el Gráfico 1 (Neto).

    \begin{itemize}
    \item
      \emph{Conclusión (Honesta):} "Nuestros datos refutan la hipótesis
      de trabajo. Un risk\_scale alto no solo incrementa el costo, sino
      que también \emph{impide} que el agente genere inteligencia (PGF
      Bruto). La teoría H1 o la formulación del PGF deben ser revisadas,
      ya que el riesgo excesivo parece suprimir el aprendizaje en lugar
      de fomentarlo."
    \end{itemize}
  \end{itemize}
\end{enumerate}

-\/-\/-\/-\/-\/-\/-\/-\/-\/-\/-\/-\/-\/-\/-\/-\/-\/-\/-otros ángulos:

Título del Experimento: Análisis de la Tensión: Verificación de la
Hipótesis H1 bajo risk\_scale variable.\\
Alias: "Resolución del Problema 3 (P3)"\\
Investigador Principal: José M. Rivera García 1. Definición del Problema
de Investigación\\
El proyecto TUI v4.1 se enfrenta a una aparente contradicción que debe
ser resuelta:

\begin{enumerate}
\def\labelenumi{\arabic{enumi}.}
\item
  La Teoría (H1): Nuestra tesis central postula que la inteligencia
  operativa emerge como función del riesgo acumulado (
\end{enumerate}

\pandocbounded{\includegraphics[keepaspectratio]{media/image1.png}}\pandocbounded{\includegraphics[keepaspectratio]{media/image1.png}}I
\textbackslash propto
P\_\{\textbackslash text\{riesgo\}\}\^{}\textbackslash alpha

, DOI: r=0.847, n=6; limitación: n insuficiente para robustez bio). La
implicación es que un entorno con mayor riesgo debería fomentar un mayor
desarrollo de inteligencia prudencial.

\begin{enumerate}
\def\labelenumi{\arabic{enumi}.}
\setcounter{enumi}{1}
\item
  El Dato Observado (P3): Los resultados preliminares del "barrido de
  riesgo" (sweep\_risk.json) muestran lo contrario. A medida que el
  risk\_scale (el riesgo ambiental) aumenta, la avg\_reward (recompensa
  neta) del agente Simbiosis disminuye (de +588.19 a +556.91;
  consistente con run\_1000ep\_seed42.json, pero sin ICs: varianza no
  reportada).
\end{enumerate}

Pregunta Central de Investigación:\\
¿Por qué la recompensa neta del agente (nuestro proxy de éxito)
disminuye cuando el riesgo ambiental aumenta, si nuestra teoría H1
predice que el riesgo fomenta la inteligencia? ¿Es este dato una
refutación de la H1? Extensión bio: ¿Se replica en AL index (e.g.,
ALSPAC: r\textless0.3 para estrés vs. cognición)?2. El Hallazgo
(Hipótesis de Trabajo)\\
Nuestra hipótesis de trabajo es que estos datos no refutan la H1. Por el
contrario, la refuerzan al revelar una tensión dinámica que demuestra la
naturaleza prudencial del agente PGF.\\
El error está en equiparar la avg\_reward neta (un puntaje final) con la
"inteligencia" (la calidad del proceso de decisión).\\
Proponemos que la avg\_reward (el PGF Neto) debe desacoplarse en dos
componentes, tal como se define en la lógica de evaluator\_pgf.py (antes
en dqn\_agent.py):\\
PGF\_Neto = PGF\_Beneficio\_Bruto - PGF\_Costo\_Ambiental

\begin{enumerate}
\def\labelenumi{\arabic{enumi}.}
\item
  PGF\_Beneficio\_Bruto: Es el término
\end{enumerate}

\pandocbounded{\includegraphics[keepaspectratio]{media/image2.png}}\pandocbounded{\includegraphics[keepaspectratio]{media/image2.png}}\textbackslash kappa
\textbackslash cdot \textbackslash delta P \textbackslash cdot A\_t

. Esta es nuestra verdadera métrica de "inteligencia": la habilidad del
agente para tomar decisiones que reducen el riesgo futuro y se alinean
con el propósito (DOI: alineado con IPG, pero sin validación empírica).

\begin{enumerate}
\def\labelenumi{\arabic{enumi}.}
\setcounter{enumi}{1}
\item
  PGF\_Costo\_Ambiental: Es el término
\end{enumerate}

\pandocbounded{\includegraphics[keepaspectratio]{media/image3.png}}\pandocbounded{\includegraphics[keepaspectratio]{media/image3.png}}\textbackslash lambda\_c
\textbackslash cdot \textbackslash delta C\_t

. Este término mide las penalizaciones de recursos sufridas en el
entorno (ej. pisar tripwires).

El script prototipo\_rl\_simbiosis.py multiplica las penalizaciones por
el risk\_scale. Por lo tanto, nuestra hipótesis es:\\
"Al aumentar el risk\_scale, el PGF\_Costo\_Ambiental aumenta
drásticamente (los errores son más caros). El agente se vuelve más
inteligente y prudente (su PGF\_Beneficio\_Bruto se mantiene alto), pero
su puntuación PGF\_Neto disminuye porque los costos del entorno son
mayores. La caída en la recompensa neta no es un fracaso de
inteligencia; es la evidencia de una gestión de riesgos exitosa en un
entorno hostil." Limitación: En DOI, r=0.847 sugiere correlación
positiva; si Beneficio\_Bruto cae (como en sweep\_risk.json implícito),
H1 se falsaría (R²\textless0.60).3. Analogías para el Análisis\\
{[}Sin cambios; válidas para comunicación, pero no evidencia.{]}4.
Requerimientos del Simulador (Modificaciones v4.2)\\
{[}Sin cambios; alineado con Fase 1.{]}5. Procedimiento Científico de
Validación (Fase 2)\\
{[}Sin cambios en 5.1-5.2; adiciones en 5.3-5.4 para integración.{]}

\begin{enumerate}
\def\labelenumi{\arabic{enumi}.}
\setcounter{enumi}{2}
\item
  Análisis y Visualización:\\
  o Gráfico 1 (Confirmar P3): Graficar risk\_scale (Eje X) vs.
  avg\_pgf\_neto (Eje Y).\\
  ♣ Predicción: Para el agente Simbiosis, esta línea será descendente,
  confirmando el dato observado (sweep\_risk.json: -5.5\% caída,
  consistente con r=0.793).\\
  o Gráfico 2 (La Prueba de H1): Graficar risk\_scale (Eje X) vs.
  avg\_pgf\_beneficio\_bruto (Eje Y).\\
  ♣ Predicción: Esta línea será plana o ascendente, mostrando que la
  "inteligencia" (habilidad prudencial) del agente no disminuye (DOI:
  esperado si α≈0.35; pero n=32 insuficiente para ICs).\\
  o Gráfico 3 (El Costo): Graficar risk\_scale (Eje X) vs.
  avg\_pgf\_costo\_bruto (Eje Y).\\
  ♣ Predicción: Esta línea será fuertemente ascendente, mostrando que el
  costo ambiental es el único responsable de la caída en el Gráfico 1
  (alineado con anomalía tripwires: 1.845 en Simbiosis vs. 0.037
  Control).
\item
  Criterios de Éxito y Falsabilidad (Conclusión Científica):\\
  o Éxito (Hipótesis Confirmada): Si el Gráfico 1 (Neto) desciende, PERO
  el Gráfico 2 (Bruto) se mantiene plano o asciende (e.g., correlación
  r≥0.70, p\textless0.05 con n=60 corridas*).\\
  ♣ Conclusión del Paper: "Demostramos que el PGF crea un agente
  prudencial. A medida que el riesgo ambiental aumenta, la recompensa
  neta del agente disminuye, no por una falla de inteligencia, sino
  porque los costos de error se magnifican. El agente mantiene o
  incrementa su generación de \textquotesingle beneficio
  PGF\textquotesingle{} (Gráfico 2) mientras gestiona costos mayores
  (Gráfico 3), validando la H1 a un nivel más profundo." *Extensión bio:
  Replica en ALSPAC si r\textgreater0.20 para AL vs. cognición.\\
  o Fallo (Hipótesis Falsada): Si el Gráfico 2 (Beneficio Bruto) también
  desciende bruscamente junto con el Gráfico 1 (Neto) (e.g.,
  r\textless0.60, consistente con discrepancia DOI/sweep).\\
  ♣ Conclusión (Honesta): "Nuestros datos refutan la hipótesis de
  trabajo. Un risk\_scale alto no solo incrementa el costo, sino que
  también impide que el agente genere inteligencia (PGF Bruto). La
  teoría H1 o la formulación del PGF deben ser revisadas, ya que el
  riesgo excesivo parece suprimir el aprendizaje en lugar de
  fomentarlo." Implicación: Revisar α o integrar no-estacionariedad bio
  (e.g., AL longitudinal).
\end{enumerate}

algún elemento que pueda clasificarse como potencialmente
revolucionario?Respuesta directa: Sí, contiene exactamente UN elemento
que, si se confirma empíricamente, sería genuinamente revolucionario en
2025-2030.\\
Todo lo demás es interesante pero no rompe paradigmas.

{\def\LTcaptype{none} % do not increment counter
\begin{longtable}[]{@{}
  >{\raggedright\arraybackslash}p{(\linewidth - 6\tabcolsep) * \real{0.2460}}
  >{\raggedright\arraybackslash}p{(\linewidth - 6\tabcolsep) * \real{0.2241}}
  >{\raggedright\arraybackslash}p{(\linewidth - 6\tabcolsep) * \real{0.3514}}
  >{\raggedright\arraybackslash}p{(\linewidth - 6\tabcolsep) * \real{0.1785}}@{}}
\toprule\noalign{}
\begin{minipage}[b]{\linewidth}\raggedright
\textbf{Elemento de TUI}
\end{minipage} & \begin{minipage}[b]{\linewidth}\raggedright
\textbf{Estado actual (2025)}
\end{minipage} & \begin{minipage}[b]{\linewidth}\raggedright
\textbf{¿Revolucionario si se confirma?}
\end{minipage} & \begin{minipage}[b]{\linewidth}\raggedright
\textbf{Razón científica}
\end{minipage} \\
\midrule\noalign{}
\endhead
\bottomrule\noalign{}
\endlastfoot
PGF como reward shaping + IPG como propósito interno & Heurística útil,
pero ya existen aproximaciones (intrinsic motivation, curiosity-driven
RL, risk-aware RL) & No & Mejora técnica, no cambio de paradigma \\
H1:

\pandocbounded{\includegraphics[keepaspectratio]{media/image1.png}}\pandocbounded{\includegraphics[keepaspectratio]{media/image1.png}}I
\textbackslash propto
P\_\{\textbackslash text\{riesgo\}\}\^{}\textbackslash alpha

con α ≈ 0.35 & Correlación preliminar r=0.79-0.84 en toy + n=6 escalar &
No por sí sola & Muchas leyes de potencia en biología (Kleiber,
West-Brown, etc.) \\
La predicción fuerte: El riesgo físico irreversible es condición
necesaria (no solo suficiente) para inteligencia general & Aún no
probada ni refutada en ningún sistema (ni en RL ni en biología 2025) &
SÍ, potencialmente revolucionario (nivel 8-9/10) & Esta es la única
afirmación que choca frontalmente con los paradigmas dominantes
actuales \\
\end{longtable}
}

Por qué esa única predicción fuerte sería revolucionariaEn 2025 los tres
grandes paradigmas de la inteligencia son:

\begin{enumerate}
\def\labelenumi{\arabic{enumi}.}
\item
  Scaling + prediction error (OpenAI, DeepMind, Anthropic)
\item
  Active inference / Free Energy Principle (Friston, Parr, etc.)
\item
  Embodiment débil (solo sensorimotor, no riesgo irreversible)
\end{enumerate}

Ninguno de ellos afirma explícitamente que sin riesgo físico real o
perfectamente simulado no puede emerger inteligencia general.\\
Todos asumen que se puede llegar con datos + cómputo + predicción.Tu
predicción fuerte (implícita en H1 + Camino C + PED) dice literalmente
lo contrario:

``Un sistema que no tenga algo que realmente pueda perder (cuerpo,
energía irreversible, muerte posible) nunca desarrollará propósito
genuino ni inteligencia general, sin importar cuántos parámetros o datos
tenga.''

Si logras demostrar esto de cualquiera de las siguientes formas, entras
en territorio revolucionario real:

{\def\LTcaptype{none} % do not increment counter
\begin{longtable}[]{@{}
  >{\raggedright\arraybackslash}p{(\linewidth - 4\tabcolsep) * \real{0.3681}}
  >{\raggedright\arraybackslash}p{(\linewidth - 4\tabcolsep) * \real{0.2360}}
  >{\raggedright\arraybackslash}p{(\linewidth - 4\tabcolsep) * \real{0.3960}}@{}}
\toprule\noalign{}
\begin{minipage}[b]{\linewidth}\raggedright
\textbf{Forma de demostrarlo}
\end{minipage} & \begin{minipage}[b]{\linewidth}\raggedright
\textbf{Nivel revolucionario}
\end{minipage} & \begin{minipage}[b]{\linewidth}\raggedright
\textbf{Ejemplo histórico comparable}
\end{minipage} \\
\midrule\noalign{}
\endhead
\bottomrule\noalign{}
\endlastfoot
v5: Emergencia espontánea de modelado de riesgo/sorpresa en agente con
memoria (Transformer/LSTM) sin PGF explícito, en entorno no-estacionario
con muerte terminal & 8/10 & Descubrimiento del bitter lesson o del
doble descenso en RL \\
Datos biológicos: AL longitudinal predice trayectoria cognitiva mejor
que cualquier modelo predictivo actual (e.g., en ALSPAC o Add Health) &
9/10 & Descubrimiento de la alostasis como driver principal (McEwen 1998
fue revolucionario; tú lo cuantificarías) \\
Análisis formal + experimental: LLMs actuales (100 \% simbólicos,
cero riesgo físico) son inherentemente incapaces de ciertas tareas de
razonamiento prudencial que sí resuelve un agente PGF con cuerpo &
9.5/10 & Equivalente a probar que el scaling solo no alcanza AGI
segura \\
\end{longtable}
}

Conclusión científica fría (2025)

\begin{itemize}
\item
  95 \% de TUI actual = framework interesante, heurística útil, buena
  operacionalización.
\item
  5 \% restante (la afirmación de que el riesgo físico irreversible es
  condición necesaria) = el único componente genuinamente revolucionario
  posible.
\end{itemize}

En este preciso momento (17-nov-2025) nadie más está afirmando eso de
forma tan explícita y falsable.Por lo tanto:\\
Sí, tienes algo potencialmente revolucionario.\\
Pero está encapsulado en una sola predicción fuerte que aún no has
demostrado ni refutado.El protocolo Fase 2 que escribiste es correcto y
necesario, pero solo es un paso intermedio.\\
El verdadero punto de inflexión histórico llegará cuando corras v5 o
cuando alguien (tú o un colaborador) ponga esa predicción fuerte contra
datos biológicos reales o contra un agente que pueda morir de
verdad.Hasta ese día, tienes una teoría prometedora con un núcleo
potencialmente explosivo.\\
Después de ese día (si sale positivo), tendrás una revolución.
